\documentclass[12pt]{article}

% --- Codificación e idioma ---
\usepackage[utf8]{inputenc}
\usepackage[T1]{fontenc}
\usepackage[spanish]{babel}

% --- Diseño de página ---
\usepackage[a4paper, left=2.54cm, right=2.54cm, top=2.54cm, bottom=3.17cm]{geometry}

% --- Paquetes útiles ---
\usepackage{graphicx}
\usepackage{tabularx}
\usepackage{booktabs}
\usepackage{pgfgantt}
\usepackage{xcolor}
\usepackage{fancyhdr}
\usepackage{lastpage}
\usepackage{hyperref}
\usepackage{float}

% --- Paquetes locales ---
\usepackage{python_style}
\usepackage{json_style}
\usepackage{c_style}

% --- Colores del Gantt ---
\definecolor{barblue}{RGB}{0,0,102}
\definecolor{barpurple}{RGB}{102,0,102}

% --- Encabezado y pie ---
\pagestyle{fancy}
\fancyhf{}
\lhead{\centering ESCUELA DE INGENIERÍA EN COMPUTACIÓN Y TELECOMUNICACIONES}
\rfoot{\thepage\ de \pageref{LastPage}}
\lfoot{Tutor Socrático -- Reporte de Status}

\begin{document}

% --- Título ---
\begin{center}
    \textbf{Tutor Socrático: Sistema Inteligente de Tutoría Basado en LLM con Metodología Socrática para la Enseñanza de Introducción a la Algoritmia} \\
    \textbf{Reporte 1}
\end{center}

\vspace{1em}

% --- Datos generales ---
\begin{table}[h]
    \begin{tabularx}{\textwidth}{|l|X|}
        \hline
        \textbf{Fecha del Reporte} & 02/03/2026 \\
        \hline
        \textbf{Integrantes del Proyecto} & Cristian Ignacio de la Hoz Reyes, 10149779 \\
                                          & Manuel Jose Rodriguez Cruz, 10150681 \\
        \hline
        \textbf{Periodo} & \\
        Desde: & 01/07/2026 \\
        \hline
        \textbf{A:} & 04/14/2026 \\
        \hline
    \end{tabularx}
\end{table}

\vspace{1em}

\newpage

% --- Rendimiento del calendario ---
\begin{table}[h]
    \begin{tabularx}{\textwidth}{|l|X|}
        \hline
        \textbf{Rendimiento del Calendario} & [Variación del Calendario] \% adelantado/atrasado \\
        \hline
        65\% & \\
        \hline
    \end{tabularx}
\end{table}

% --- Gantt Chart ---
\begin{center}
    \resizebox{\textwidth}{!}{
        \begin{ganttchart}[
            canvas/.style={draw=none},
            hgrid,
            vgrid,
            x unit=1.5cm,
            y unit title=0.7cm,
            y unit chart=0.6cm,
            title height=1,
            bar/.style={fill=barblue, draw=none},
            bar height=0.7,
            group right shift=0,
            group top shift=.6,
            bar label font=\small\color{black!80},
            title label font=\small\bfseries
        ]{1}{7}
        
            \gantttitle{Timeline}{7} \\
            \gantttitle{26-Ene}{1}
            \gantttitle{12-Feb}{1}
            \gantttitle{26-Feb}{1}
            \gantttitle{2-Mar}{1}
            \gantttitle{17-Mar}{1}
            \gantttitle{23-Mar}{1} 
            \gantttitle{30-Mar}{1} \\
        
            \ganttbar{Arquitectura del Proyecto}{1}{1}
            \ganttset{bar/.style={fill=none}}
            \ganttbar[inline, bar label font=\footnotesize]{Done}{2}{2} \\
            \ganttset{bar/.style={fill=barblue, draw=none}}
        
            \ganttbar{Extracción materiales curriculares}{1}{2}
            \ganttset{bar/.style={fill=none}}
            \ganttbar[inline, bar label font=\footnotesize]{Done}{3}{3} \\
            \ganttset{bar/.style={fill=barblue, draw=none}}
        
            \ganttbar{Construccion dialogos socraticos}{2}{4}
            \ganttset{bar/.style={fill=none}}
            \ganttbar[inline, bar label font=\footnotesize]{Done}{5}{5} \\
            \ganttset{bar/.style={fill=barblue, draw=none}}
          
            \ganttbar{Chunking y Vectorización}{2}{4} 
            \ganttset{bar/.style={fill=none}}
            \ganttbar[inline, bar label font=\footnotesize]{Done}{5}{5} \\
            \ganttset{bar/.style={fill=barblue, draw=none}}
            
            \ganttbar{Arquitectura RAG}{3}{4} 
            \ganttset{bar/.style={fill=none}}
            \ganttbar[inline, bar label font=\footnotesize]{Done}{5}{5} \\
            \ganttset{bar/.style={fill=barblue, draw=none}}
            \ganttbar{Guardrails de seguridad y pedagogía}{3}{5} \\
            \ganttbar{Interfaz Conversacional inicial}{6}{6} \\
            \ganttbar[bar/.style={fill=barpurple}]{Buffer Zone}{6}{7}
            
        \end{ganttchart}
    }
\end{center}

\vspace{1em}

% --- Problemas actuales ---
\begin{table}[h]
    \begin{tabularx}{\textwidth}{|l|X|}
        \hline
        \multicolumn{2}{|l|}{\textbf{Mayores problemas existentes}} \\
        \hline
        1 & Recaudación de data previamente estructurada para la transformación requerida para el módulo de RAG. \\
        \hline
        2 & Recaudación y creación de ejemplos socráticos basados en fuentes propietarias de alta calidad. \\
        \hline
    \end{tabularx}
\end{table}

\vspace{1em}

% --- Acciones planificadas ---
\begin{table}[h]
    \begin{tabularx}{\textwidth}{|l|X|}
        \hline
        \multicolumn{2}{|l|}{\textbf{Acciones planificadas para resolver los problemas existentes}} \\
        \hline
        1 & Extender la búsqueda con técnicas avanzadas en motores de búsqueda basadas en fuentes de alta calidad. \\
        \hline
        2 & Aumentar reuniones con entidades pertinentes para creación de documentos del syllabus del proyecto. \\
        \hline
    \end{tabularx}
\end{table}

\vspace{1em}

% --- Temas que requieren atención ---
\begin{table}[h]
    \begin{tabularx}{\textwidth}{|l|X|}
        \hline
        \multicolumn{2}{|l|}{\textbf{Temas que requieren la atención}} \\
        \hline
        1 & Una máquina dedicada con recursos necesarios para desarrollo y funcionamiento adecuado del proyecto. \\
        \hline
    \end{tabularx}
\end{table}

\newpage

% --- Logros del período ---
\begin{table}[h]
    \begin{tabularx}{\textwidth}{|l|X|}
        \hline
        \multicolumn{2}{|l|}{\textbf{Mayores logros del periodo}} \\
        \hline
        1 & Recopilación y estructuración del corpus de diálogos socráticos reales facilitados por el Prof. Alfredo Vásquez. \\
        \hline
        2 & Generación de 85 conversaciones sintéticas preservando la metodología y voz del profesor, totalizando 88 conversaciones de entrenamiento con 1.824 turnos. \\
        \hline
        3 & Implementación del pipeline de extracción, segmentación y vectorización del corpus de conocimiento: 2.725 fragmentos producidos desde 4 fuentes heterogéneas. \\
        \hline
        4 & Ingesta exitosa de 2.603 vectores en ChromaDB con recuperación semántica validada (similitud coseno $\approx 0.74$--$0.85$). \\
        \hline
    \end{tabularx}
\end{table}

\vspace{1em}

% ===========================================================================
\section*{Descripción del Progreso}
% ===========================================================================

El presente reporte documenta el avance correspondiente al primer objetivo
específico del proyecto: la recopilación, estructuración y vectorización del
corpus de conocimiento pedagógico del programa ICC-101-T. A continuación se
describe cada componente implementado, acompañado de ejemplos representativos
del trabajo realizado.

% ---------------------------------------------------------------------------
\subsection*{1. Recopilación del corpus pedagógico real}
% ---------------------------------------------------------------------------

El punto de partida del sistema fue la obtención de material auténtico generado
en el contexto real de la asignatura. El profesor Alfredo Vásquez facilitó tres
transcripciones de sesiones de tutoría socrática cubiertas las cuales contemplan
los temas de bucles \texttt{while}, funciones y bucles anidados. Estos diálogos
constituyen la semilla metodológica del sistema: de ellos se extrae la voz
del docente, el arco pedagógico que va de la analogía del mundo real hacia la
construcción incremental del código, y el registro informal en español que
caracteriza las sesiones presenciales de la asignatura.

Cada transcripción fue normalizada a un formato multi-turno JSON con tres roles
diferenciados: \texttt{system} (instrucción de comportamiento del tutor),
\texttt{user} (respuestas del estudiante) y \texttt{assistant} (intervenciones
del profesor). El fragmento a continuación ilustra los primeros turnos del
diálogo sobre bucles \texttt{while}, en el cual el profesor construye la
intuición del problema antes de introducir una sola línea de código:

\begin{jsoncode}
{
  "source_file": "synthetic",
  "topic": "bucle while — monto total de N ventas",
  "conversations": [
    {
      "role": "system",
      "content": "Eres un tutor socrático de programación en lenguaje C.
Tu objetivo no es dar respuestas directas, sino guiar al estudiante
mediante preguntas que lo lleven a descubrir el concepto por sí mismo.
Usa ejemplos del mundo real antes de introducir código."
    },
    {
      "role": "user",
      "content": "Profe, tengo que hacer un programa que calcule el monto
total de N ventas. ¿Cómo lo hago con un while?"
    },
    {
      "role": "assistant",
      "content": "Antes de escribir una sola línea de código, pensemos en
el problema como si fuéramos el cajero de una tienda. Imagínate que estás
en la caja registradora y van llegando clientes uno por uno. El primero
compra por 150, el segundo por 230, el tercero por 80. ¿Qué harías tú,
con lápiz y papel, para saber cuánto vendiste en total al final del día?"
    },
    {
      "role": "user",
      "content": "Iría anotando cada venta y al final las sumaría todas."
    },
    {
      "role": "assistant",
      "content": "Bien. ¿Y cómo las sumarías? ¿Esperas a tener todas y las
sumas de un golpe, o vas acumulando a medida que llega cada cliente?"
    },
    {
      "role": "user",
      "content": "Voy acumulando profe, cada vez que llega uno le sumo su
monto al total que llevo."
    },
    {
      "role": "assistant",
      "content": "Exactamente. Entonces tienes una variable mental que va
creciendo con cada venta. ¿Con qué valor empieza esa variable antes de
que llegue el primer cliente?"
    }
  ]
}
\end{jsoncode}

Nótese el patrón socrático: el tutor no enuncia la definición de acumulador
directamente. En su lugar, construye la abstracción a través de una situación
cotidiana — la caja registradora — llevando al estudiante a articular por sí
mismo el concepto de variable acumuladora y la necesidad de inicializarla en
cero antes de entrar al ciclo. Esta secuencia didáctica, repetida de forma
consistente a lo largo de todos los diálogos, es precisamente el patrón que el
sistema debe aprender e internalizar.

% ---------------------------------------------------------------------------
\subsection*{2. Generación del corpus sintético}
% ---------------------------------------------------------------------------

Dado que tres diálogos resultan insuficientes para el ajuste fino de un modelo
de lenguaje, se procedió a escalar el corpus mediante la generación de 85
conversaciones sintéticas que preservan la voz, el arco pedagógico y el
registro informal del profesor, extendiendo la cobertura a la totalidad de los
temas del programa ICC-101-T contemplados por el proyecto.

La estrategia de variación operó sobre tres ejes simultáneos: el tema
algorítmico (desde bucles hasta recursión y arreglos multidimensionales), el
dominio del problema del mundo real utilizado como analogía introductoria
(restaurantes, hospitales, torneos deportivos, fábricas, aeropuertos, entre
otros), y el tipo de error conceptual del estudiante que el tutor debe
redirigir sin revelar la respuesta. Cada diálogo generado fue sometido a una
lista de verificación que garantiza la ausencia de respuestas directas
tempranas, la presencia de al menos un error del estudiante con redirección
socrática, y el cierre con un ejercicio propuesto.

El dataset consolidado en \texttt{data/dataset.jsonl} presenta la siguiente
distribución:

\begin{table}[H]
    \centering
    \begin{tabular}{lrr}
        \toprule
        \textbf{Categoría} & \textbf{Conversaciones} & \textbf{Turnos aprox.} \\
        \midrule
        Bucle \texttt{while}       & 15 & 310 \\
        Bucle \texttt{for}         & 10 & 207 \\
        Bucle \texttt{do-while}    &  8 & 165 \\
        Bucles anidados            & 10 & 207 \\
        Funciones                  & 10 & 207 \\
        Arreglos 1D                & 10 & 207 \\
        Arreglos 2D                &  8 & 165 \\
        Estructuras (\textit{structs}) & 8 & 165 \\
        Recursión                  &  6 & 124 \\
        Diálogos reales (semilla)  &  3 &  67 \\
        \midrule
        \textbf{Total}             & \textbf{88} & \textbf{1.824} \\
        \bottomrule
    \end{tabular}
    \caption{Distribución del dataset de entrenamiento por categoría temática.}
\end{table}

El fragmento siguiente corresponde al diálogo \texttt{arrays1d\_01\_notas.json}
y ejemplifica el tratamiento de un tema no cubierto por los diálogos reales. La
analogía de la lista de calificaciones en papel conduce al estudiante a
descubrir la necesidad de los arreglos frente a la limitación de las variables
individuales, y el concepto de índice base cero emerge de la misma interacción
antes de ver cualquier sintaxis de C:

\begin{jsoncode}
{
  "topic": "Arreglos 1D — guardar N notas y mostrarlas al final",
  "conversations": [
    {
      "role": "assistant",
      "content": "Antes de pensar en código, imagina que eres un profesor
corrigiendo exámenes. Tienes 30 exámenes en la mesa y necesitas anotar
la nota de cada uno para leerlas al final en voz alta. ¿Cómo lo harías
en papel?"
    },
    {
      "role": "user",
      "content": "Las escribiría en una lista, una debajo de la otra."
    },
    {
      "role": "assistant",
      "content": "Exactoo. Una lista numerada. Ahora dime: si en C quisieras
guardar solo 3 notas, ¿qué variables declararías?"
    },
    {
      "role": "user",
      "content": "int nota1, nota2, nota3;"
    },
    {
      "role": "assistant",
      "content": "¿Y si el grupo puede tener entre 1 y 200 estudiantes y
no sabes el número hasta que el programa corre? ¿Cuántas variables
declararías?"
    },
    {
      "role": "user",
      "content": "Mmm… no sé, ¿declaro 200 variables?"
    },
    {
      "role": "assistant",
      "content": "Piénsalo un momento. Si el grupo tiene 5 estudiantes,
desperdiciarías 195. Si mañana tiene 201, el programa falla. ¿Crees que
declarar variables individuales funciona cuando N es desconocido en
tiempo de compilación?"
    }
  ]
}
\end{jsoncode}

% ---------------------------------------------------------------------------
\subsection*{3. Pipeline de extracción, segmentación y vectorización (RAG)}
% ---------------------------------------------------------------------------

En paralelo con la construcción del corpus de entrenamiento, se implementó el
pipeline de Generación Aumentada por Recuperación (\textit{Retrieval-Augmented
Generation}, RAG) que proveerá al tutor de contexto técnico preciso durante las
sesiones de tutoría. La arquitectura general del sistema se ilustra en la
figura siguiente.

\begin{figure}[H]
    \centering
    \includegraphics[width=1\textwidth]{Tesis RAG.png}
    \caption{Arquitectura del pipeline RAG del Tutor Socrático.}
\end{figure}

El pipeline, implementado en el notebook \texttt{01-chuncking.ipynb}, opera en
tres etapas sucesivas: ingestión de documentos crudos, segmentación semántica,
y agrupación en fragmentos (\textit{chunks}) incrustables. Se definieron tres
abstracciones de datos centrales y cuatro ingestores especializados:

\begin{python}
@dataclass
class RawDocument:
    """Salida de cualquier ingestor."""
    doc_id:      str          # id del capítulo o URL
    source_name: str          # tutorialspoint | litmentor | epub | web
    source_type: SourceType
    content:     str          # HTML crudo o Markdown
    title:       str = ""
    metadata:    dict = field(default_factory=dict)

@dataclass
class Segment:
    """Unidad atómica de parsing."""
    doc_id:          str
    source_name:     str
    title:           str
    segment_type:    SegmentType   # text | code | table | list
    content:         str
    heading_context: str = ""
    token_count:     int = 0

@dataclass
class Chunk:
    """Unidad final incrustable."""
    id:       str    # chunk_00001
    source:   str    # fuente de origen
    url:      str    # URL o id de capítulo EPUB
    title:    str    # título de página/capítulo
    heading:  str    # sección dentro de la página
    content:  str    # markdown limpio (texto embebido y enviado al LLM)
    has_code: bool
    tokens:   int

    @property
    def content_hash(self) -> str:
        return hashlib.md5(self.content.encode()).hexdigest()
\end{python}

Las cuatro fuentes de conocimiento procesadas corresponden a: el libro
\textit{Modern C} en formato EPUB (22 capítulos), un conjunto de 148 páginas
de Tutorialspoint cubiertas mediante scraping por lotes, 48 páginas de
LitMentor, y el sitio \texttt{how2lab.com} recorrido en su totalidad mediante
rastreo automático. Para el scraping web se empleó una instancia local de
Firecrawl, lo cual permite extraer contenido limpio en Markdown desde páginas
con renderizado dinámico sin depender de servicios externos. Los resultados de
la ejecución del pipeline se resumen a continuación:

\begin{table}[H]
    \centering
    \begin{tabular}{lrrrr}
        \toprule
        \textbf{Fuente} & \textbf{Docs} & \textbf{Segmentos} & \textbf{Chunks} & \textbf{Con código} \\
        \midrule
        \textit{Modern C} (EPUB)    &  22 &  3.258 &  922 & 548 (59\%) \\
        TutorialsPoint              & 148 & 31.869 &  882 & 570 (65\%) \\
        LitMentor                   &  48 &  1.366 &  366 & 283 (77\%) \\
        how2lab.com (crawl)         &  55 &  5.794 &  555 & 416 (75\%) \\
        \midrule
        \textbf{Total}              & \textbf{273} & \textbf{42.287} & \textbf{2.725} & \textbf{1.817 (67\%)} \\
        \bottomrule
    \end{tabular}
    \caption{Resultados del pipeline de extracción y segmentación por fuente.}
\end{table}

El chunker semántico (\texttt{SemanticChunker}) opera con un presupuesto de
2.048 tokens por fragmento utilizando el tokenizador \texttt{cl100k\_base} de
tiktoken. Los segmentos de código actúan como anclas: el chunker agrega texto
circundante siempre que el presupuesto lo permita, garantizando que cada chunk
incluya su contexto explicativo. Un control adicional verifica que el total
de tokens con el prefijo de incrustación \texttt{search\_document:} no exceda
el límite absoluto de 8.192 tokens del modelo \texttt{nomic-embed-text}:

\begin{python}
MAX_CHUNK_TOKENS = 2048
_NOMIC_HARD_LIMIT = 8192
_EMBED_PREFIX = "search_document: "
_EMBED_PREFIX_TOKENS = len(TOKENIZER.encode(_EMBED_PREFIX))

# Límite seguro = límite del modelo - tokens del prefijo - margen de 50
EMBED_MAX_TOKENS = _NOMIC_HARD_LIMIT - _EMBED_PREFIX_TOKENS - 50  # = 8139

def count_embed_tokens(text: str) -> int:
    """Cuenta tokens tal como nomic-embed-text los verá (prefijo incluido)."""
    return _EMBED_PREFIX_TOKENS + count_tokens(text)

def exceeds_embed_limit(text: str) -> bool:
    return count_embed_tokens(text) > EMBED_MAX_TOKENS
\end{python}

% ---------------------------------------------------------------------------
\subsection*{4. Ingesta vectorial y validación de recuperación}
% ---------------------------------------------------------------------------

Los 2.725 fragmentos exportados en \texttt{chunks\_all.json} fueron procesados
en el notebook \texttt{02\_chromadb\_ingestion.ipynb}. Tras filtrar 88 chunks
de contenido insuficiente (menos de 100 caracteres), se procedió a la
incrustación mediante \texttt{nomic-embed-text} ejecutado localmente a través
de Ollama. El protocolo de prefijos distingue entre documentos y consultas,
respetando el esquema de incrustación asimétrica propio del modelo:

\begin{python}
EMBED_MODEL     = "nomic-embed-text"
COLLECTION_NAME = "socratic_tutor_collection"
BATCH_SIZE      = 50

def get_embedding(text: str, is_query: bool = False) -> list:
    prefix = "search_query: " if is_query else "search_document: "
    response: EmbeddingsResponse = ollama.embeddings(
        model=EMBED_MODEL,
        prompt=prefix + text,
    )
    return response.embeddings

def embed_with_retry(text: str, is_query: bool = False) -> list:
    """Incrustación con backoff exponencial para errores transitorios."""
    last_err = None
    for attempt in range(MAX_RETRIES + 1):
        try:
            return get_embedding(text, is_query=is_query)
        except Exception as e:
            if "input length exceeds" in str(e).lower():
                raise  # No reintentar chunks sobredimensionados
            last_err = e
            if attempt < MAX_RETRIES:
                wait = RETRY_DELAY * (2 ** attempt)
                time.sleep(wait)
    raise last_err

# Almacenamiento en ChromaDB (distancia coseno)
collection.add(
    ids=[c["id"]],
    embeddings=[embedding],
    documents=[c["content"]],
    metadatas=[{"id": c["id"], "heading": c["heading"],
                "has_code": str(c["has_code"])}],
)
\end{python}

De los 2.637 chunks candidatos, 2.603 fueron ingestados exitosamente en
aproximadamente 161 segundos (16,5 chunks/s). Los 34 chunks restantes excedieron
el límite del modelo en tiempo de ejecución a pesar de haber pasado el
preflight basado en tiktoken; esto se debe a una discrepancia entre el
tokenizador \texttt{cl100k\_base} y el tokenizador interno de Ollama para
textos muy densos en caracteres especiales, y constituye un punto de mejora
identificado para el próximo período.

La validación de la colección mediante consultas semánticas sobre temas
centrales del syllabus confirmó que el motor de recuperación funciona
correctamente. Las similitudes coseno obtenidas se muestran a continuación:

\begin{table}[H]
    \centering
    \begin{tabular}{lr}
        \toprule
        \textbf{Consulta de validación} & \textbf{Similitud coseno} \\
        \midrule
        \textit{What is a struct in C?}              & 0.85 \\
        \textit{How do arrays relate to pointers?}   & 0.81 \\
        \textit{What are function pointers?}         & 0.84 \\
        \textit{How do pointers work in C?}          & 0.81 \\
        \textit{How does a for loop work in C?}      & 0.79 \\
        \textit{How do you handle strings in C?}     & 0.75 \\
        \textit{What is undefined behavior in C?}    & 0.74 \\
        \textit{Difference between malloc and calloc?} & 0.71 \\
        \bottomrule
    \end{tabular}
    \caption{Similitudes coseno obtenidas en las consultas de validación sobre la colección ChromaDB.}
\end{table}

Los resultados confirman que el índice vectorial recupera fragmentos altamente
relevantes para consultas sobre los temas del programa, con similitudes que
oscilan entre 0.71 y 0.85 en el rango de consultas evaluadas.

\begin{figure}[H]
    \centering
    % [PLACEHOLDER: insertar aquí una captura de pantalla del output de las
    %  celdas de validación del notebook 02_chromadb_ingestion.ipynb,
    %  mostrando los resultados de las consultas con sus similitudes,
    %  encabezados de los chunks recuperados y fragmentos de código.]
    \fbox{\parbox{0.85\textwidth}{\centering\vspace{2cm}
        \textit{[Captura: resultados de consultas de validación en ChromaDB ---
        notebook \texttt{02\_chromadb\_ingestion.ipynb}]}
    \vspace{2cm}}}
    \caption{Salida representativa del módulo de validación de recuperación semántica.}
\end{figure}

\vspace{1em}

% --- Logros próximos ---
\begin{table}[H]
    \begin{tabularx}{\textwidth}{|l|X|}
        \hline
        \multicolumn{2}{|l|}{\textbf{Logros planificados para el período siguiente}} \\
        \hline
        1 & Implementación de una interfaz conversacional inicial que integre el módulo RAG con un LLM base. \\
        \hline
        2 & Diseño e implementación de guardrails pedagógicos y de seguridad para acotar las respuestas del tutor al dominio de la asignatura. \\
        \hline
        3 & Inicio del ajuste fino (\textit{fine-tuning}) del modelo base con el dataset de diálogos socráticos generado, utilizando el framework Unsloth con adaptadores LoRA. \\
        \hline
    \end{tabularx}
\end{table}

\end{document}
